\section{Methodology}
\label{sec:methodology}

\subsection{Parameter space}
\label{sec:parameter-space}

The first step of the study was to identify the parameters that could
affect the performance of an IO500 execution. The configuration of IO500
contains parameters related to the I/O benchmarks, but the performance
also depends on the execution environment, including the compute nodes
and the parallel file system. We therefore considered parameters from
three main aspects: the IO500 benchmarks, the compute environment, and the
storage system.


Table~\ref{tab:parameter-space} summarises the parameters identified in
these different parts of the system. The role of each group of parameters
is explained in more detail in the following paragraphs.

\begin{table}[H]
\centering
\caption{Parameters considered in the IO500 tuning study.}
\label{tab:parameter-space}

\rowcolors{2}{gray!10}{white}
\begin{tabular}{lll}
    \hline
    \rowcolor{gray!30}
    \textbf{Category} & \textbf{Parameter} & \textbf{Description} \\
    \hline

    
    IOR & \texttt{transferSize} & Size of each I/O operation \\
    IOR & \texttt{blockSize} & Data handled by each process \\
    IOR & \texttt{filePerProc} & One file per process \\
    
    IOR & \texttt{segmentCount} & Number of file segments \\
    \hline

    MDTEST & \texttt{n} & Number of files per process \\
    MDTEST & \texttt{files-per-dir} & Number of files per directory \\
    \hline

    FIND & \texttt{nproc} & Number of processes \\
    FIND & \texttt{pfind-queue-length} & Size of the work queue \\
    FIND & \texttt{pfind-parallelize...} & Parallel directory access \\
    \hline

    Machine & \texttt{nodes} & Number of compute nodes \\
    Machine & \texttt{processes-per-node} & Number of processes per node \\
    \hline

    BeeGFS & \texttt{num-targets} & Number of storage targets used \\
    \hline
\end{tabular}


\end{table}



\paragraph{IOR parameters.}

The IOR parameters mainly control the way data is transferred during the
bandwidth benchmarks. The \texttt{transferSize} determines the amount of
data transferred by each I/O operation. Changing this value can modify the
number of I/O requests sent to the storage system and therefore affect the
measured bandwidth.

The \texttt{blockSize} determines the amount of data handled by each
process. Together, \texttt{transferSize} and \texttt{blockSize} define an
important part of the I/O access pattern generated by the benchmark.

The \texttt{filePerProc} parameter determines whether each process works
on its own file or whether several processes access a shared file. This
choice can strongly change the level of contention and the way the
parallel file system handles the benchmark.

Finally, \texttt{segmentCount} controls the number of segments used by the
IOR benchmark. Increasing the number of segments changes the total amount
of data processed and the sequence of I/O operations performed by each
process.

\paragraph{MDTEST parameters.}

The MDTEST parameters control the metadata benchmark generated by IO500.
The parameter \texttt{n} determines the number of files handled by each
process. A larger value therefore generates more file-system metadata
operations.

The \texttt{files-per-dir} parameter controls how these files are
distributed between directories. This can affect performance because
placing many files in the same directory does not generate the same
metadata access pattern as distributing them across several directories.

\paragraph{Find parameters.}

The FIND benchmark evaluates the performance of file discovery operations.
The \texttt{nproc} parameter determines the number of processes involved
in the search and therefore controls the level of parallelism.

The \texttt{pfind-queue-length} parameter controls the amount of work that
can be kept in the search queue. The
\texttt{pfind-parallelize-single-dir-access-using-hashing} parameter
allows accesses to a directory to be distributed between several
processes. These parameters influence how the search work is divided and
executed in parallel.

\paragraph{Compute environment parameters.}

The number of compute nodes and the number of processes per node also
affect IO500 performance. The \texttt{nodes} parameter determines how many
\texttt{bora} nodes participate in an execution, while
\texttt{processes-per-node} determines how many benchmark processes are
started on each node.

Together, these two parameters determine the total number of processes
accessing the storage system. Increasing this number can provide more I/O
parallelism, but it can also increase contention when many processes
access the storage system simultaneously.

\paragraph{Parallel file system parameter.}

The experiments use BeeGFS as the parallel file system. The
\texttt{num-targets} parameter controls the number of storage targets used
to distribute file data.

Using a larger number of targets can distribute I/O operations across
more storage resources and potentially increase the available bandwidth.
However, the best value depends on the benchmark and on the number of
processes accessing the storage system. The effect of this parameter must
therefore be considered together with the IOR and compute parameters.

\paragraph{From the parameters to the search space.}

These parameters do not act independently. For example, increasing the
number of processes can change the benefit obtained from using more
storage targets, while the effect of a given \texttt{transferSize} can
depend on the level of parallelism used during the execution.

A configuration is therefore defined by a combination of values from the
different parameter categories. The set of all possible combinations
forms the parameter space of the study. Exploring this space makes it
possible to study the interactions between parameters and to identify
configurations that provide high IO500 scores.

However, the number of possible combinations quickly becomes large. This
makes a complete evaluation expensive when each configuration requires an
IO500 execution. The following sections describe how this parameter space
was explored and how a reduced-cost exhaustive campaign was used as a
reference for the evaluation of the search strategies.


\subsection{Exhaustive search}
\label{sec:exhaustive-search}

\subsubsection{Search space}
\label{sec:search-space}

\subsubsection{Cheap evaluation setup}
\label{sec:cheap-evaluation}

\subsubsection{Repetitions}
\label{sec:exhaustive-repetitions}

\subsection{Ground-truth construction}
\label{sec:ground-truth}

\subsection{Search strategies}
\label{sec:search-strategies}

\subsubsection{Random search}
\label{sec:random-search}

\subsubsection{Bayesian Optimization}
\label{sec:bo-methodology}

\subsubsection{BO configurations}
\label{sec:bo-configurations}

% \subsection{Experimental protocol}
% \label{sec:experimental-protocol}

% \subsubsection{Evaluation budget}
% \label{sec:evaluation-budget}

% \subsubsection{Random seeds}
% \label{sec:random-seeds}

% \subsubsection{Repetitions}
% \label{sec:protocol-repetitions}

% \subsubsection{Offline evaluation}
% \label{sec:offline-evaluation}

% \subsection{Validation campaign}
% \label{sec:validation}

% \subsection{Human specialist baseline}
% \label{sec:human-baseline}